\chapter{Políticas de Visitas y TI, pruebas de seguridad y evidencias}

Implementación de clientes, automatización y matriz: Genser Andree Catalán Espina. Ejecución de
pruebas y recopilación de evidencias: Pablo Daniel Barillas Moreno y Andrés Rafael Chivalán. Este
capítulo describe los clientes temporales que representan a los segmentos de Visitas y TI, las
reglas que implementan sus políticas y la metodología con la que se comprueba toda la matriz de
tráfico de la Fase 2. Su estado es de
\emph{implementación local validada}; los resultados en AWS se incorporan cuando existan salidas
y capturas reales de la ventana autorizada.

\section{Clientes de prueba}

\begin{table}[H]
\centering
\caption{Clientes temporales de la Fase 2}
\small
\begin{tabularx}{\textwidth}{>{\raggedright\arraybackslash}p{2.4cm}p{2.3cm}p{3.6cm}Y}
\toprule
\textbf{Cliente} & \textbf{IP privada} & \textbf{IPv4 pública} & \textbf{Security Group} \\
\midrule
Visitas & \code{10.10.0.4} & Efímera, solo durante la ventana & \code{visitas\_client}: sin entradas desde Internet \\
TI & \code{10.10.0.132} & No & \code{ti}: única entrada SSH desde el bastión \\
\bottomrule
\end{tabularx}
\end{table}

Ambos clientes se definen en \code{infra/terraform/test-clients.tf}, dependen de
\code{enable\_phase2\_compute} (falso por defecto), usan Amazon Linux 2023 con IMDSv2 obligatorio
y disco EBS cifrado, y se etiquetan como temporales. La subred de Visitas no asigna IPv4 pública
por defecto, por lo que la salida a Internet se declara de forma explícita con
\code{associate\_public\_ip\_address = true}; la dirección desaparece al destruir la instancia.
El \code{user\_data} (\code{templates/test-client-user-data.tftpl}) instala la consulta DNS
\code{rednova-dnsquery}, instala y verifica \code{nftables} y ejecuta el hardening común. TI
descarga el paquete mediante el Gateway Endpoint de S3, sin salida general a Internet. Además,
espera una respuesta válida de BIND9, configura \code{systemd-resolved}, verifica la resolución
con \code{getent} y usa una copia controlada de \code{/etc/resolv.conf} como respaldo si el modo
de resolución de Amazon Linux no aplica el dominio de enrutamiento.

\section{Implementación de las políticas}

\begin{enumerate}
  \item \textbf{Visitas sin acceso interno}: la NACL de Visitas deniega la VPC en la regla 100
    antes de permitir HTTP y HTTPS hacia Internet en las reglas 110 y 120. Solo se exceptúan el
    resolvedor de la VPC (\code{10.10.0.2}) y las respuestas a TI y al bastión. Su Security
    Group, además, solo permite salida DNS al resolvedor y TCP 80/443.
  \item \textbf{TI hacia todas las subredes}: el Security Group de TI permite iniciar cualquier
    tráfico hacia la VPC. Para que las respuestas ICMP regresen, los Security Groups de Ventas,
    bastión, DNS, intranet y web aceptan ICMP \emph{echo-request} desde \code{10.10.0.128/27}, y
    las NACL de Visitas y DMZ (salida) y de TI (entrada) admiten únicamente ICMP
    \emph{echo-reply} (tipo 0).
  \item \textbf{Nadie inicia conexiones hacia TI}: la NACL de TI solo acepta SSH desde
    Administración, puertos efímeros de retorno y \emph{echo-reply}; todo lo demás cae en la
    regla 200. Un puerto alto que superara la NACL se detiene en el Security Group de TI, que
    no tiene entradas salvo SSH desde el bastión.
\end{enumerate}

\begin{infobox}[Defecto detectado antes del despliegue]
El validador estático mostró que, con las reglas heredadas, TI no podía hacer ping a ninguna
subred: las NACL descartaban el \emph{echo-reply} y los Security Groups de destino no aceptaban
ICMP. Se corrigió con los cambios mínimos descritos en el punto 2, sin abrir ninguna entrada
hacia TI.
\end{infobox}

\section{Metodología de pruebas}

La suite \code{scripts/test-phase2.ps1} ejecuta cada prueba desde su segmento de origen por SSH
a través del bastión, o desde el equipo del operador para las pruebas desde Internet. Las sondas
(\code{scripts/phase2-probe.sh}, \code{dnsquery.py} y \code{check-public-ports.py}) se envían por
la entrada estándar, de modo que no se instala nada en las instancias. El script:

\begin{itemize}
  \item exige indicar quién autorizó la ventana y pide confirmación antes de generar tráfico;
  \item no ejecuta Terraform ni modifica Security Groups, NACL o instancias;
  \item distingue un bloqueo de red (tiempo agotado) de un destino alcanzado (respuesta o RST)
    y de un error de configuración (SSH, DNS o ruta);
  \item guarda por prueba fecha, commit, origen, destino, comando, salida, esperado y obtenido,
    con identificadores de cuenta redactados;
  \item termina con código 1 si algún resultado contradice la matriz.
\end{itemize}

Pablo opera la suite y organiza las salidas por identificador. Andrés Rafael confirma la identidad,
región, estado y permisos de la ventana, participa en la ejecución y revisa que las evidencias no
expongan secretos. Cada archivo registra a ambos como responsables operativos; la autoría técnica
de los clientes y scripts permanece atribuida a Genser.

Antes de la ventana, \code{scripts/validate-phase2.py} evalúa 50 flujos contra los Security
Groups y NACL declarados en Terraform, capa por capa e incluyendo el tráfico de retorno.

\begin{table}[H]
\centering
\caption{Pruebas de Visitas, TI y seguridad}
\small
\begin{tabularx}{\textwidth}{>{\raggedright\arraybackslash}p{1.4cm}p{2.5cm}p{5.2cm}Y}
\toprule
\textbf{ID} & \textbf{Origen} & \textbf{Prueba} & \textbf{Esperado} \\
\midrule
F2-03 & TI & DNS y HTTP hacia \code{rrhh.rednova.internal} & Resuelve y HTTP 200 \\
F2-05 & Visitas y DMZ & HTTP, DNS y SSH hacia la intranet & Bloqueado \\
F2-06 & Internet & Resolver la intranet en el DNS público & Sin resolución \\
F2-07 & Visitas & DNS público, HTTP y HTTPS hacia Internet & Permitido \\
F2-08 & TI & ICMP y servicios en todas las subredes & Permitido \\
F2-09 & Ventas, Visitas, DMZ, Data Center & ICMP, SSH y TCP alto hacia TI & Bloqueado \\
F2-11 & Internet & SSH a la web y al cliente de Visitas & Bloqueado \\
F2-12 & Internet y TI & Escaneo de 24 puertos comunes & Solo los puertos previstos \\
\bottomrule
\end{tabularx}
\end{table}

\begin{warningbox}[Estado de las evidencias]
La validación local no sustituye a las pruebas reales. Las tablas de resultados y las capturas
\code{09} a \code{12} y \code{14} se agregan únicamente con la salida de
\code{test-phase2.ps1} obtenida en la ventana AWS autorizada.
\end{warningbox}
